\section{Watched Literals}

Before we get into the Watched Literal algorithm, I want to ease us into it.
Why? Because it's so different from our current version of unit propgation. 
For a moment (and for the rest of the book) let's forget about pure literal 
elimination; it wasn't that great anyways. Unit propagation gave us a massive speedup 
oppossed to recursive backtracking, taking runtimes down from minutes to seconds and 
from half-hours to mere minutes. Despite all of that, it's really inefficient in the 
sense that we're going so much work for so little in return.

Let's define what a \textbf{contradiction} is, because we've been using that word 
here and there, but beyond ``unsatisfiable,'' we haven't formally defined what it means.
A \textbf{contradiction} is the result of a literal forcing a clause to evaluate to False,
leading to the whole assignment to become unsatisfiable. You might remember how in the 
brute force solver, we added an \textit{if()} statement at the end of evaluating a singular 
clause that checked if that clause evaluated to False. If it did, we skipped that assignment 
and moved onto the next one, but if it didn't, we continue our evaluation. 

This led to a pretty significant speedup in the brute forcer because instead of evaluating everything 
then deciding, we did it incrementally, allowing us to quickly leave when something was bad. 
At the moment, our original idea for unit propagation follows that first idea of delaying 
evaluation to the end. However, if you go back and look at the pseudocode, you'll notice that we 
do something \textit{almost} the same as checking if a clause evaluates to False, at the moment,
we check if everything except for one element evaluates to False, and that one element is Unassigned
(since we made it so that unit propagation only operates on unsatisfied clauses). However, notice 
that when a clause evaluates to False, rather than stopping our propagation and assignment, we 
continue until we reach the evaluation phase? 

Just like in brute force, the idea is the same: once a clause evaluates to False, no matter 
where we are in our assignment, there's no point in continuing on because we've found out that 
nothing will unfalsify that clause. As an exercise, I want you to try adding that into your 
unit propagation code and seeing how fast it takes to run on our test sets. I would do it, but 
I have a \underline{massive} headache from implementing this algorithm last night.
\footnote{This headache was also induced by my lack of ability to fall asleep at all last night
AND a really intense session at the gym. Computer science gives you headaches for sure, but not like 
this one.}
Once you're done, we'll discuss watched literals some more.

So that early contradiction detection is one of the ideas behind watched literals. 
It's not one the key ideas, but it is part of it. There are some other ideas behind 
watched literals. For instance: zero cost backtracking. Not O(1) backtracking, $0$
backtracking. Speaking of cost, do you notice something else wasteful about the 
way we propagate? If not for our \textit{satisfied} marker array, we'd have to 
rescan the whole equation every time we propagated, which is about twice per level. 
That's pretty wasteful. After all, why do have to scan every clause in their entirety
every single time? Especially because, usually, clauses don't become unsatisfied very 
often. And because they're not at risk of becoming unsatisfied very often, why do we need 
to scan it? And besides, it's not like every single variable appears in every single 
clause, so why not scan only the clauses we have to? I hope these questions are leading 
you to think of other questions regarding inefficiencies in our current implementation.

These questions were the driving force into the invention of the 
\textit{Watched Literals} algorithm.

\textbf{History can be found here}
\footnote{\url{https://school.a4cp.org/summer2011/slides/Gent/SATCP3.pdf}}

Let's learn how the Watched Literals algorithm works. Earlier, I made a comment 
about how usually, each clause doesn't contain every single variable in the 
equation. So when we force a variable's assignment, we don't need to check 
every single clause, only at most the ones that contain the variable and its 
negation. (Ie, if we force a variable, say $2 = T$, we only need to check the clauses 
containing $2$ and the clauses containing $-2$). However, do we \textit{need} to 
check the clauses containing $2$? Think about it. Those clauses are already 
satisfied by $2 = T$, so there's nothing to propagate on, for those clauses. 
The only clauses that have a \textit{chance} to produce a propagation are the 
ones containing $-2$ because $-2 = F$, so if they have one Unassigned variable 
and everything else is now False, we can propagate!

\subsection{A Visual Example of Watched Literals}
Let's view a visual example. Suppose we had a clause with 5 literals 
(doesn't matter what they are). The first two literals are watched
\begin{center}
    \[ \blackinwhitesquare \blackinwhitesquare \square \square \square \]
    \[ U U U U U \]
\end{center}
Suppose the first literal now evaluates to False. We now scan the clause for 
another literal that doesn't evaluate to False. We will find it at the third literal.
\begin{center}
    \[ \square \blackinwhitesquare \blackinwhitesquare \square \square \]
    \[ F U U U U \]
\end{center}
Now suppose the fifth literal becomes False. We do nothing because we're not watching 
the fifth literal.
\begin{center}
    \[ \square \blackinwhitesquare \blackinwhitesquare \square \square \]
    \[ F U U U F \]
\end{center}
Now the third literal becomes False. We scan again for another not-False literal.
\begin{center}
    \[ \square \blackinwhitesquare \square \blackinwhitesquare \square \]
    \[ F U F U F \]
\end{center}
Now the second literal becomes False. Other than the second watched literal (the fourth),
there are no more not-False literals. We don't relocate the watch.
\begin{center}
    \[ \square \blackinwhitesquare \square \blackinwhitesquare \square \]
    \[ F F F U F \]
\end{center}
However, now the second watched literal knows it's the last Unassigned literal in the clause.
It will now evaluate to True!
\begin{center}
    \[ \square \blackinwhitesquare \square \blackinwhitesquare \square \]
    \[ F F F T F \]
\end{center}
Our clause is now satisfied. Suppose instead, when we were here:
\begin{center}
    \[ \square \blackinwhitesquare \square \blackinwhitesquare \square \]
    \[ F F F U F \]
\end{center}
That the fourth literal also became False.
\begin{center}
    \[ \square \blackinwhitesquare \square \blackinwhitesquare \square \]
    \[ F F F F F \]
\end{center}
We've now encountered a \textit{contradiction}, meaning that whatever assignment we have
is impossible to satisfy the equation. Suppose we backtrack, and that restores us to a 
state where two literals in this same clause are now Unassigned. What happens to the 
watched literals?
\begin{center}
    \[ \square \square \square \square \square \]
    \[ U U F F F \]
\end{center}
The answer is that they stay where they are.
\begin{center}
    \[ \square \blackinwhitesquare \square \blackinwhitesquare \square \]
    \[ U U F F F \]
\end{center}
This is because one watched literal is Unassigned, so everything is okay. Although 
this last diagram isn't accurate to what a real restored state would look like, we can
suspend our disbelief for the sake of simplicitly. The reason why watched literals work is 
because of the \textbf{Watched Literal Invariant}, which is that at any given point in time,
at least one of the watched literals is either True or Unassigned. This invariant is what
allowed us to propagate the clause when there was only one Unassigned literal left, and we 
will see further how it helps.


\subsection{Algorithmic Changes}
\begin{lstlisting}
public interface:
    SATSolver(string input)
    bool solveCNF()

private interface:
    enum var
    {
        False
        True
        Unassigned
    }

    void parse(string equation)
    bool solve(var[] assignment, int level)
    optional int[] delta propagate(var[] assignment)
    bool createWatched()

    int numVars
    int numClauses

    int[][] clauses
    bool[] satisfied
    queue<int> unitProp

    optional var[] vars
    map<int, int[]> var_to_clause
    map<int, (int, int)> clause_to_vars
\end{lstlisting}

So firstly, there are some major differences. We removed the \textit{evaluate()} method. 
This is because now that propagation checks for contradictions, having all variables 
assigned is now \textbf{\textit{equivalent}} to finding a satisfactory assignment of 
variables. So we don't need to evaluate our assignment anymore! That's significantly 
less work than what we had before, which is what we want. Next, we added in a queue, 
\textit{unitProp}, kind of like what we had for unit propagation before. This time,
it's now a member of the solver object instead of a local object within the method.
We will get back to that in a second. Lastly, we added in two maps, one from 
variables to clauses, and the other from clauses to a pair of variables. 

These data structures are used for fast lookup between variables to  
\textbf{\textit{indices}} of clauses (for \textit{var\_to\_clause}), and 
one from \textbf{\textit{indices}} of clauses to \textbf{\textit{indices}} of 
watched literals. It's important that we remember we're using indices instead 
of actual clauses because forgetting how we're using our mapping is going to create 
a lot of bugs and we're gonna have a bad time. 

We also added in a new method, called \textit{createWatched()} and what it does 
is it creates our initial maps for var\_to\_clause and clause\_to\_vars. 
Here is the pseudocode:

\begin{algorithm}[H]
\caption{Watched Literals Creation}\label{alg:16}
\begin{algorithmic}[1]

\Procedure{createWatched}{}

\For{each $Clause$ in $Clauses$}
    \If{$Clause$ has at least 2 literals}
        \State lit1 $\gets$ Clause[0]
        \State lits $\gets$ Clause[1]
        \State cIdx $\gets$ $\Call{indexOf}{Clauses, Clause}$
        \State $\Call{append}{varToClause[lit1], cIdx}$
        \State $\Call{append}{varToClause[lit2], cIdx}$
        \State $clauseToVar[cIdx] \gets (0, 1)$
    \ElsIf{$Clause$ is has 1 literal}
        \State lit $\gets$ Clause[0]
        \State idx $\gets \Call{abs}{lit}$ 
        \State cIdx $\gets$ $\Call{indexOf}{Clauses, Clause}$
        \State $\Call{append}{varToClause[lit], cIdx}$
        \State $clauseToVar[cIdx] \gets (0, 0)$
        \State $\Call{enqueue}{unitProp, var}$

        \If{lit > 0} 
            \State $assignment[idx] \gets true$
        \Else
            \State $assignment[idx] \gets false$
        \EndIf
    \Else
        \State return $false$
    \EndIf
\EndFor

\State return $true$

\EndProcedure
\end{algorithmic}
\end{algorithm}

\noindent So here's how to read it: for each clause, it will fall into one of three cases:
\begin{enumerate}
    \item It's size is greater than or equal to 2
    \item It only has 1 variable
    \item It's empty
\end{enumerate}

In the first case, the two literals that will be watching that clause will be the 
first two literals. So the index $i$ gets appended to their array of clause indices.
Then, we record that clause $i$ is being watched by its first and second literal.

In the second case, since there's only one literal, it gets watched by only that 
literal (duh). We also record that clause $i$ is being watched by only its first literal.
And lastly, we enqueue that literal to our unit propagation queue so we can do a root 
level unit propagation before we start solving. We also want to immediately force the 
assignment of the variable so that way we can use its value immediately in solving.

In the last case, since there are no literals, this clause is automatically unsatisfiable, 
so therefore this whole equation is unsatisfiable, so we return False so that when we call 
this method, we know that we can't solve it. 

Now let's look at our modified solving methods:

\begin{algorithm}[H]
\caption{Modified Setup}\label{alg:17}
\begin{algorithmic}[1]

\Procedure{solveCNF}{}

\State assignment[1..numVars]

\For{each $var$ in $assignment$}
    \State $var \gets Unassigned$
\EndFor

\If{\Call{createWatched}{} = false}
    \State return $false$
\ElsIf{\Call{propagate}{assignment, 0} = false}
    \State return $false$
\EndIf

\State return $\Call{solve}{assignment, 1}$

\EndProcedure
\end{algorithmic}
\end{algorithm}

There are some slight differences: we create our watched literals and make sure 
that there isn't an unsatisfiable clause, and we make sure that \textit{propagate()}
returns something. If both those things are successful, then we can start solving away.

Next, we'll modify our \textit{solve()} method slightly:

\begin{algorithm}[H]
\caption{Modified solving method}\label{alg:18}
\begin{algorithmic}[1]
\Procedure{solve}{assignment[1..n], level}

\If{level = numVars}
    \State return $true$
\EndIf

\If{assignment[level] $\neq$ Unassigned}
    \State return $\Call{solve}{assignment, level + 1}$
\EndIf

\State assignment[level] $\gets true$
\State pResult $\gets \Call{propagate}{assignment, level}$

\If{propagation failed}
    \State assignment[level] $\gets false$
    \State pResult $\gets \Call{propagate}{assignment, -level}$

    \If{propagation failed}
        \State return $false$
    \EndIf

    \If{\Call{solve}{assignment, level + 1} = $true$}
        \State return $true$
    \EndIf

    \Comment solving failed, so undo propagation changes

    \State \Call{undoPropagation}{pResult}
    \State return $false$
\EndIf

\If{\Call{solve}{assignment, level + 1} = $true$}
    \State return $true$
\EndIf

\State \Call{undoPropagation}{pResult}

\State assignment[level] $\gets false$
\State pResult $\gets \Call{propagate}{assignment, -level}$

\If{propagation failed}
    \State return $false$
\EndIf

\If{\Call{solve}{assignment, level + 1} = $true$}
    \State return $true$
\EndIf

\State \Call{undoPropagation}{pResult}
\State assignment[level] $\gets Unassigned$
\State return $false$

\EndProcedure
\end{algorithmic}
\end{algorithm}

This is almost the same as our old unit propagation code, except for that we don't 
keep track of which clauses are satisfied across levels, and the propagation method 
returns an array of variable indices instead of the actual variables stored in \textit{pResult}.
Undoing propagation involves unassigning variables involved in the propagation and marking 
the newly satisfied clauses as unsatisfied. The code is also split up into multiple branches 
depending on whether or not propagation was successful. Some parts of the code seem really 
similar to each other because of this. I am sure that there may be a way to simplify the similar 
parts, but I like the way I wrote it because I find it readable and the first time I wrote it, it 
was very nested.

Now for the interesting part:

\begin{algorithm}[H]
\caption{Watched Literal Propagation}\label{alg:19}
\begin{algorithmic}[1]
\Procedure{propagate}{assignment[1..n], int decision}

\State delta[1..m]

\If{decision $\neq$ 0}
    \State $\Call{enqueue}{unitProp, decision}$
    \State idx $\gets \Call{abs}{decision}$
    \State $\Call{append}{delta, idx}$
\EndIf

\While{unitProp is not empty}
    \State prop $\gets$ \Call{dequeue}{unitProp}

    \If{varToClause does not contain -prop}
        \State continue on to the next loop interation
    \EndIf

    \State watchedClauses $\gets$ varToClause[prop]

    \For{each $cIdx$ in watchedClauses}
        \State Clause[1..t] $\gets$ Clauses[cIdx]
        \State watches $\gets$ clauseToVar[cIdx] \Comment pair of clause indices

        \If{Clause[watches.first] $\neq$ -prop}
            \State $\Call{swap}{watches.first, watches.second}$
        \EndIf

        \State relocated $\gets false$

        \For{i $\gets$ 0 up to \Call{size}{Clause}}
            \If{i = watches.first $\vee$ i = watches.second}
                \State continue
            \EndIf

            \State unwatched $\gets$ Clause[i]
            \State uIdx $\gets \Call{abs}{unwatched}$

            \If{\Call{evalsFalse}{unwatched} = $false$}
                \State $\Call{append}{varToClause[unwatched], cIdx}$
                \State \Comment erase the clause index at its index within watchedClauses
                \State wIdx $\gets \Call{indexOf}{watchedClauses, cIdx}$ 
                \State $\Call{eraseAt}{watchedClauses, wIdx}$
                \State relocated $\gets true$
                \State break
            \EndIf
        \EndFor

        \If{relocated = $false$}
            \If {$\Call{checkValid}{Clause, watches.second, delta}$ = null}
                \State return null
            \EndIf
        \EndIf
    \EndFor
\EndWhile

\State return delta

\EndProcedure
\end{algorithmic}
\end{algorithm}

\begin{algorithm}[H]
\caption{check if the current assignment is valid}\label{alg:20}
\begin{algorithmic}[1]
\Procedure{checkValid}{Clause[1..n], secondWatch, delta[1..m]}

\State otherWatch $\gets$ Clause[secondWatch]
\State oIdx $\gets \Call{abs}{otherWatch}$

\If{assignment[oIdx] = Unassigned}
    \If{otherWatch > 0}
        \State assignment[oIdx] $\gets true$
    \Else 
        \State assignment[oIdx] $\gets false$
    \EndIf

    \State $\Call{append}{delta}{oIdx}$
    \State $\Call{enqueue}{unitProp}{otherWatch}$
\Else
    \If{\Call{evalsFalse}{otherWatch} = false}
        \State continue onto the next loop iteration
    \Else 
        \State $\Call{undoPropagation}{delta}$
        \State $\Call{clear}{unitProp}$
        \State return null
    \EndIf
\EndIf

\EndProcedure
\end{algorithmic}
\end{algorithm}

This is a huge method, (so large it had to be split into two pieces so it would fit)
so let's walkthrough it so we understand every aspect of it. If we look back at our \textit{solve()} 
code, we pass in \textit{cur\_var} and -\textit{cur\_var} as the \textit{decision} parameter into
the propagation method. And in the beginning of \textit{solveCNF()}, we pass in $0$, to signify propagation at the 
root level. When that happens, we don't add anything into the queue, instead we just propagate 
on the literals that were forced during the creation of the watched literals. 

Like I said before, we only need to check the clauses containing the negative version of the 
literal being propagated on because those are the clauses that the actual unit propagation could 
be forced upon. However, we only need to actually check the clauses being watched by the negative 
version of the literal being propagated on. 

Next, we grab the indices of those clauses being watched and then for each one, we try 
to relocate the index of the first watched literal to a new literal that doesn't evaluate 
to False. Why only the first? Because swapping the two indices and operating on only the 
first index is far easier than creating a second branch dedicated to the second index.
If a literal doesn't evaluate to False, we first make that unwatched literal now watch 
the current clause, we then make the current watched literal stop watching the current clause,
and then we update the first index of the watched literals to be the index of the newly 
watching literal. That means we have successfully relocated the watched literal and no 
unit propagation needs to occur. Because we're modifying the array of clause indices, 
we don't want to increment the index variable in our implementation.

\textit{Aside:
When you delete something in a dynamic array, in order to fill the blank space, every element to 
the right of what you deleted has to shift over to the left by one unit. This takes O(n) time and 
on large arrays can be very costly. Because the next element has shifted over to the left, its index 
is now the current index, so you don't want to increment the index variable. 
In the implementation, the relocation logic is implemented as a ``swap and pop'' in which we swap the 
current element with the last one, then we delete the last element, which was the current element. 
This leads to no shifting of elements, and is logically equivalent (for our use case), even if the 
order of elements, relative to the deleted element, is no longer the same.
}

If we couldn't relocate the watched literal's index, we want to increment the index variable 
so that on the next iteration it will point to a new element. In the case that we couldn't 
relocate the watched literal, it means that every element in the clause, except for possibly the 
element pointed to by the second watched literal, evaluates to False. If the element pointed to by 
the second watched literal also evaluates to False, then we've encountered a contradiction in our 
assignment. However, if the second watched literal evaluates to True, then the clause is already 
satisfied and we need to do nothing except for move onto the next variable to propagate on. 
And if the second watched literal evaluates to Unassigned, then we can assign it the value it 
needs to evaluate to True and we then propagate on that newly assigned unit literal. 

If we encounter a contradiction, we must undo all the propagation we've done, including 
the initial decision before returning a null value to signify that the propagation failed. 
However, if it has succeeded, then we return the array of the variable indices that we've 
assigned.

So one of the big aspects of watched literals is how we get away with doing so little 
amount of work. And the answer is something called the ``watched literal invariant.''
This invariant is that at all times, at least one of the watched literals doesn't 
evaluate to False. I don't expect this to make everything click immediately, so let's 
explain more. A clause is only a contradiction when everything evaluates to False.
However, we don't want to check every single clause when we assign only one variable.
And even checking the negative of the variable we assigned is too much. So instead, we 
change which variables are watching which clauses every now and then. We pick at most 
two variables to watch a singular clause with, and each one watches their respective clauses
until they are assigned a value that makes them evaluate to False. In that case, there is 
the possibility that the watched literal invariant may be violated (ie, there is a possibility
that both watched literals may evaluate to False), but just to be sure, we try to change the 
variable watching that same clause. 

We scan the clause for any unwatched literal that is Unassigned or True, and if we find one, 
we've restored the watched literal invariant, so we don't need to propagate or do anything 
else for that matter. However, if we can't find one, then we're in serious danger of the 
watched literal invariant being violated. So our only hope is for the second watched literal 
to either already be True (we need to do nothing), or Unassigned (we assign it and propagate).

Thanks to that invariant, we also don't need to do any undoing when we backtrack other 
than undoing what we propagated. So long as the watched literal invariant holds true 
when we undo our assignments, the watched literals don't need to move, which means 
besides the undoing of what we propagated, we get zero cost backtracking.

\subsection{Why Two Watches?}
This now begs the question, why two watched literals? What is special about the number $2$?
Why can't we have just one watched literal? In order to answer this question, we need to ask 
ourselves: ``what can we do with two watched literals that we can't do with one?'' 
When we have two watched literals, we can efficiently check for whether or not we can 
propagate. How so?

After we try propagating on a literal, we scan the whole clause except for the indices 
of the first and second watched literals. Next, if everything in the clause evaluates to 
False, our only hope to propagate is on the second watched literal. So if that evaluates to 
False as well, we have found a contradiction, and if it doesn't we have found a literal we 
can propagate on.

But why is that more efficient that just one watched literal? Let's think of the worst case 
scenario for if we just had one watched literal. What does that look like? Suppose we were 
watching the first literal in the clause and every other literal in the clause evaluated 
to False. Ie, we need to propagate on this watched literal. How would it know it needs to 
propagate when it was never alerted by anything becoming False? It would have no idea 
it needs to do anything until that watched literal was falsified, and by that point,
everything in the clause would evaluate to False, so the clause would be unsatisfied,
only then would we backtrack, but when we backtrack, we'd still have no way of knowing if 
the clause is unit. 

So with only one watched literal, in the worst case scenario, which is that the literal
we're watching is the last literal in the clause to be falsified, we'd be doing the same 
amount of work as normal unit propagation, perhaps even more. 

Visually speaking, this is what I'm talking about:
\begin{center}
    \[ \blackinwhitesquare \square \square \]
    \[ U U U \]
\end{center}
The last literal becomes False, but we're not watching it so nothing moves.
\begin{center}
    \[ \blackinwhitesquare \square \square \]
    \[ U U F \]
\end{center}
Next the second literal becomes False, but we're not watching it so nothing happens.
\begin{center}
    \[ \blackinwhitesquare \square \square \]
    \[ U F F \]
\end{center}
Although we should perform unit propagation, our watched literal is completely blind 
to the fact that it is the last Unassigned literal in the clause, so it remains as is.
If luck has it, it could become True and all would be well, but this is reality, so more 
likely than not, it will also become False.
\begin{center}
    \[ \blackinwhitesquare \square \square \]
    \[ F F F \]
\end{center}
Now we scan the clause, and we see that we can't relocate the watched literal, so we 
backtrack on our assignment.
\begin{center}
    \[ \blackinwhitesquare \square \square \]
    \[ U F F \]
\end{center}
Since the watched literal was never aware if both of the other literals became False 
in one pass of the algorithm or if it was across multiple decisions, it still doesn't 
know that it's the last Unassigned literal. It has to scan the clause each time to know 
whether or not it's the last Unassigned literal, which is no better than regular unit
propagation.

\subsection{Watched Literals - Results}

After implementing watched literals, our runtime decreases dramatically. 
For uf20.91, solving all 1000 instances takes \textbf{349ms}, as opposed to the 
\textbf{1548ms} from pure literal elimination and the \textbf{868ms} from normal unit propagation.
And for UUF50.218.1000, it takes only \textbf{7130ms} as opposed to the \textbf{231,107ms} 
from pure literal elimination and the \textbf{140,411ms} from normal unit propagation.
Watched literals took us down from minutes to seconds for the unsatisfiable 
problems, which is amazing for us.

\begin{FVerbatim}
On all 1000 equations in uf20-91:

Watched Literals:          349ms (~0.35 seconds).
Pure Literal Elimination:  1548ms (~1.55 seconds).
Unit Propagation:          868ms (~0.87 seconds).
Recursive Backtracking:    1,540,613ms (~25.68 minutes).
Brute Force:               1,770,851ms (~29.51 minutes).
Short-Circuited (Basic):   49,874ms (~50 seconds).
Parallel:                  10,589ms (~10.5 seconds).

On all 1000 equations in UUF50.218.1000:

Watched Literals:         7,130ms (~7.13 seconds).
Pure Literal Elimination: 231,107ms (~3.85 minutes).
Unit Propagation:         140,041ms (~2.34 minutes).
\end{FVerbatim}